\section{Kết luận}

\subsection{Tổng kết kết quả đạt được}

Báo cáo Bài tập lớn môn Kho dữ liệu và Hệ hỗ trợ quyết định (CO4031) đã trình bày toàn bộ quá trình nghiên cứu, phân tích nghiệp vụ, thiết kế kiến trúc và hiện thực hóa hệ thống Kho dữ liệu tích hợp Hệ hỗ trợ quyết định cho bài toán doanh nghiệp bán lẻ toàn cầu Global Superstore. Đối chiếu với các mục tiêu đề ra ban đầu, nhóm thực hiện đã hoàn thành trọn vẹn toàn bộ các hạng mục kỹ thuật và nghiệp vụ từ tầng dữ liệu thô đến tầng giao diện ra quyết định.

Thứ nhất, nhóm đã tiến hành khảo sát sâu bộ dữ liệu 51.290 bản ghi giao dịch trên 147 quốc gia trong giai đoạn 2011--2014, xác lập hệ thống 13 chỉ số đo lường hiệu quả kinh doanh (KPIs) chuẩn hóa và chỉ ra các vấn đề cốt lõi về chất lượng dữ liệu nguồn cũng như bài toán thất thoát lợi nhuận do chính sách chiết khấu tại khu vực EMEA và tiểu ngành hàng Bàn làm việc (\texttt{Tables}).

Thứ hai, nhóm đã thiết kế và triển khai thành công Kho dữ liệu theo kiến trúc Medallion ba tầng (\texttt{bronze}, \texttt{silver}, \texttt{gold}) trên hệ quản trị cơ sở dữ liệu PostgreSQL 17 kết hợp Docker. Quy trình ETL được tự động hóa hoàn toàn, tích hợp cơ chế kiểm toán mã băm SHA-256, xử lý triệt để đa định dạng ngày tháng và dữ liệu mã bưu chính khuyết thiếu, đảm bảo tính lũy đẳng và hỗ trợ nạp gia tăng (Incremental Loading) thông qua cơ chế bắt thay đổi dữ liệu bằng cờ trạng thái \texttt{transform}. Mô hình hình sao (Star Schema) tại tầng Gold với 1 bảng sự kiện (\texttt{gold.fact\_sales}) và 5 bảng chiều sử dụng khóa thay thế số nguyên đã vượt qua 100\% kịch bản kiểm định chất lượng và đối soát số liệu tài chính tuyệt đối ($12.642.501,91\text{ USD}$ doanh thu, $1.467.457,29\text{ USD}$ lợi nhuận). Hệ thống 7 truy vấn phân tích OLAP và 3 bảng Materialized Views (\texttt{mv\_sales\_mart}, \texttt{mv\_product\_mart}, \texttt{mv\_customer\_mart}) đạt thời gian phản hồi thực tế dưới $50\text{ ms}$.

Thứ ba, nhóm đã phát triển hoàn chỉnh Hệ hỗ trợ quyết định (DSS) gồm ba mô-đun định lượng kết nối trực tiếp với kho dữ liệu: mô-đun dự báo chuỗi thời gian bằng thuật toán San bằng mũ có xu thế dự báo doanh thu 6 tháng đầu năm 2015 kèm khoảng tin cậy 95\% (đạt $\text{MAPE} = 14,72\%$ trên tập kiểm thử); mô-đun khai phá luật kết hợp Apriori hai cấp độ trên 25.035 giỏ hàng phát hiện 14 luật cấp tiểu ngành hàng và 27 luật cấp sản phẩm SKU (với hệ số Lift lên tới $315,57$); và mô-đun mô phỏng kịch bản What-If dựa trên độ co giãn của cầu theo giá. Toàn bộ kết quả phân tích được tích hợp đồng bộ lên dịch vụ RESTful API (FastAPI), ứng dụng Web tương tác 4 phân trang (Streamlit) và bộ báo cáo quản trị 3 trang trên Power BI Desktop, tạo thành một giải pháp hỗ trợ ra quyết định hoàn chỉnh, trực quan và có giá trị ứng dụng thực tiễn cao.

\subsection{Định hướng phát triển}

Trong giai đoạn tiếp theo, hệ thống có tiềm năng mở rộng mạnh mẽ theo hướng kiến trúc dữ liệu hiện đại quy mô doanh nghiệp lớn. Về mặt hạ tầng và kỹ thuật dữ liệu, việc tích hợp nền tảng điều phối luồng công việc tự động (Apache Airflow) và cơ chế bắt thay đổi dữ liệu dựa trên nhật ký giao dịch (Log-based CDC với Debezium và Apache Kafka) sẽ giúp chuyển đổi quy trình nạp dữ liệu từ xử lý theo lô định kỳ sang xử lý gần thời gian thực (Near Real-time Data Warehousing). Bên cạnh đó, việc nâng cấp các bảng chiều khách hàng và sản phẩm lên chuẩn SCD Type 2 hoàn chỉnh và bổ sung thêm các bảng sự kiện quản lý tồn kho (Inventory Snapshot Fact Table) sẽ giúp mở rộng phạm vi phân tích sang toàn bộ chuỗi cung ứng.

Về mặt mô hình hỗ trợ quyết định và trí tuệ nhân tạo, hệ thống có thể được nâng cấp bằng cách huấn luyện các mô hình dự báo chuỗi thời gian phân cấp (Hierarchical Time Series Forecasting) kết hợp các thuật toán học máy như Prophet, SARIMA và LightGBM để dự báo nhu cầu riêng biệt cho từng thị trường và từng nhóm ngành hàng. Đồng thời, sự kết hợp giữa ma trận phân khúc khách hàng RFM đã có sẵn trong \texttt{gold.mv\_customer\_mart} với thuật toán lọc cộng tác (Collaborative Filtering) và các luật kết hợp Apriori sẽ tạo tiền đề xây dựng một bộ máy khuyến nghị sản phẩm cá nhân hóa toàn diện, góp phần tối ưu hóa đồng thời doanh thu bán chéo và biên lợi nhuận ròng cho doanh nghiệp.
